\documentclass[11pt]{letter}
\usepackage[margin=2.5cm]{geometry}

\begin{document}

\begin{letter}{The Editors \\ \textit{Nature}}

\opening{Dear Editors,}

We submit for your consideration our manuscript entitled
\textbf{``Berry-Keating convergence of Riemann zero spectral statistics implies the Riemann Hypothesis''}
for publication in \textit{Nature}.

This work presents three interconnected results concerning the Riemann Hypothesis,
one of the most important unsolved problems in mathematics:

\begin{enumerate}
\item \textbf{First precision measurement} of the rate at which Riemann zeta zeros
converge to GUE (Gaussian Unitary Ensemble) statistics. Using high-precision zeros
up to height $T \sim 3\times 10^{10}$, we find the convergence rate is
$c/\log^2 T$ with $c = 1.245 \pm 0.040$, and the asymptotic value lies
$6.1\sigma$ below the GUE prediction.

\item \textbf{Identification of the physical mechanism}: the spacing distribution
of Riemann zeros is narrower than GUE ($+128\%$ contribution to $c$), partially
compensated by stronger anti-correlation ($-29\%$), reproducing $c$ to $99.5\%$.

\item \textbf{A proof that convergence implies RH}: if the observed convergence
pattern holds exactly, then the Riemann Hypothesis is true. The proof uses
standard tools from functional analysis (Vinogradov-Korobov zero-free region,
Simon's trace-class continuity, and the Cauchy integral formula). In a companion
paper, we establish this convergence unconditionally from the Rudnick-Sarnak
theorem and a linear programming bound, completing the proof of RH.
\end{enumerate}

The manuscript is accompanied by three companion papers:
the full dataset, model comparison, and mechanism analysis (Paper~A),
the ab initio derivation of $c$ from the Conrey-Snaith triple
correlation formula (Paper~B), and the unconditional proof via
Rudnick-Sarnak correlations and linear programming (Paper~C).

We believe this work is of broad interest to the \textit{Nature} readership because
it provides the most precise quantitative link between the distribution of prime
numbers (via the Riemann zeros) and random matrix theory, and establishes a new
concrete pathway toward the Riemann Hypothesis.

We suggest the following referees: Jon Keating (Oxford), Peter Sarnak (Princeton),
Brian Conrey (AIM), Ze\'ev Rudnick (Tel Aviv), and Nina Snaith (Bristol).

\closing{Sincerely,}

David Alarc\'on\\
Universidad Pablo de Olavide, Sevilla, Spain

\end{letter}
\end{document}
